%% The manuscript body, shared by every binding of this paper.
%% The class, the front matter and the bibliography style belong to the
%% binding; everything a reader reads belongs here, once.
\section{Introduction}

Projects that do not reach their intended evaluation tend to disappear, and the
ones that survive as text usually survive as a single sentence: the experiment
was not feasible. That sentence is almost always an aggregate over obstacles
that have nothing to do with each other. A licence that has not been granted, a
disease that is too rare in the available corpus, and a scoring instrument that
behaves unexpectedly are three different problems with three different owners,
and collapsing them into one verdict throws away the only part a reader could
act on.

This paper takes one such project and refuses the aggregate. The project set out
to evaluate chest-radiograph report generation on a public benchmark, did not
get there, and produced instead a small labelled pilot on a corpus it already
held. We report the obstacles separately, and we report what the pilot can and
cannot support, with the arithmetic that decides which is which.

The first obstacle is administrative and is not a scientific finding at all,
which is exactly why it should be stated in a scientific document. The intended
benchmark is contact-gated, the pixels underneath it carry their own data use
agreement, and both of those are ordinary and defensible choices for clinical
data. What is worth recording is their consequence: \NLocks{} obstacles were
enumerated, none of them is removed by computation, and the enumeration was
written on \BlockerDate{}, before any hardware was rented. The alternative --- to
rent a machine, discover the gate, and write the discovery up as a cost --- is
common enough to be worth naming as the thing this project avoided.

The second obstacle is statistical, and the useful move is to split it in two.
A pilot can be too small for its corpus, and a corpus can be too small for the
question. Of the \NFindings{} labelled observations, \NExceedSplit{} would need
more reports than the entire frozen test split of \NTest{} before a draw is even
expected to contain \MinPresent{} positives. For those, labelling more of the
same split changes nothing. Reporting the pilot as underpowered would be true
and would also be the wrong diagnosis.

The other half of the statistical obstacle is the draw itself, and it is the
part we did not expect. The subsample the project pre-registered is not
representative of the split it came from: it over-represents one observation at
an exact two-sided hypergeometric probability of \SkewP{}, which clears a
Bonferroni threshold of \Bonferroni{} across all \NFindings{} observations.
Because the presence gate is a threshold on positives in the draw, that anomaly
is not a nuisance detail; it is what makes one of the two scorable cells
scorable at all. Pre-registration did its job, which is to prevent choosing a
sample after seeing its labels, and it does not do the other job people expect
of it.

The third obstacle is the instrument, and it generalises furthest. The majority
control emits one constant report, the most frequent findings text in the
training split, which reads as an unremarkable study. The automatic labeller
that defines the standard metric in this area marks that text positive for
\MajAsserted{} on all \NCohort{} rows and never marks it as the observation
denoting nothing to report. A baseline that everyone treats as the floor of
predicting normality is therefore not that baseline; it is whatever the labeller
makes of one fixed string. Scored on surface text instead, the same control
reaches \SurfaceMajorityF{} token F1 against \SurfaceRetrievalF{} for a
retrieval control, a comparison that looks orderly and carries none of the
information above.

\section{Related work}

The three obstacles this paper separates have each been studied, but largely in
isolation from one another and almost always from the position of a study that
finished.

The measurement obstacle is the best documented. Automatic report generation is
conventionally scored with metrics borrowed from natural language generation,
and the most careful audit of those metrics finds that they align poorly with
what radiologists actually count as an error, which is why that work proposes
alternatives built on clinical entities and relations rather than on surface
overlap~\cite{yu2023evaluating}. That finding is the reason this paper reports
label-space counts and surface-text scores as separate, non-comparable
quantities rather than reconciling them, and the reason it does not convert its
label counts into the field's headline aggregate. It is also why the scoring
here uses the benchmark's own released labeller~\cite{smit2020chexbert} rather
than a reimplementation: if the metric is already the weak link, substituting an
approximation of it is not a neutral simplification. Our contribution on this
axis is narrow and negative --- we show that on this subsample a constant
baseline is read by that labeller as positive for a finding on every case, which
is a property of the baseline and labeller together rather than of either alone.

The power obstacle has a large literature outside this field, most influentially
the argument that low-powered studies are not merely inconclusive but actively
unreliable, because low power inflates the effect sizes of whatever does reach
significance~\cite{button2013power}. That argument is usually deployed against
studies that reported a result anyway. The situation here is the one that
argument implies but rarely addresses directly: a pilot whose subsample is too
small for most of the label schema, where the honest output is a per-cell
statement of what could and could not be scored. Separating the two power
questions --- whether the corpus could ever support a cell, and whether the
particular draw resembled the corpus --- is what keeps a corpus-level ceiling
from being misread as a sampling accident, and neither question is answered by
reporting the cells that happened to clear a gate.

The reproducibility obstacle is where the gap is widest. Community programs have
converged on code submission, artifact checklists and reproduction
challenges~\cite{pineau2021improving}, and this manuscript is built to satisfy
that standard: every number is emitted from digest-bound artifacts, and the
build fails rather than print a stale one. But such programs are designed around
a study that produced a result and asks whether someone else can obtain it
again. They have very little to say about work that could not run, and a
checklist item asking which experiments were blocked and why does not appear on
them. The consequence is that access failures are systematically invisible: they
are neither publishable results nor reproducibility defects, so they are
absorbed as private cost and the next project pays it again.

That access barriers exist in this field is uncontroversial. Credentialed access
is the norm for the major chest-radiograph
corpora~\cite{johnson2019mimiccxr,goldberger2000physionet}, and the benchmark
this project could not reach is distributed under a contact
gate~\cite{chexbench}. What is missing is not awareness of the barriers but any
convention for reporting a specific encounter with one --- which permissions
were held, which were not, what each would and would not unlock, and what the
attempt cost. This paper is an attempt to write that report for one case, and to
do it in a form that can be checked rather than asserted.

\section{Methods}
%% Independent methods file. Not woven into main.tex prose.
%% Every quantity described here resolves to a hash-bound evidence entry.

\subsection{The corpus and the frozen split}

The locally held corpus is the Indiana University chest X-ray collection
distributed through Open-i~\cite{demnerfushman2016preparing}, which pairs
de-identified radiographs with their free-text reports. Reports, not images, are
the unit of everything reported here.

The split was frozen before any of the analyses in this paper. It is
patient-disjoint by construction: a report's assignment is a function of a
salted hash of its patient identifier, so no patient contributes to more than one
side and the assignment can be recomputed by anyone holding the same identifier
list. The frozen inventory holds \NPatients{} patients and \NImages{} images,
and divides into \NTrain{} training, \NVal{} validation and \NTest{} test
reports.

One property of that inventory is worth stating because it is the kind of thing
usually smoothed over. A total of \NOutsideInventory{} radiograph files on disk
are not in
the inventory the split was built from. The project's recorded decision was to
leave them visible and outside rather than manufacture inventory rows to absorb
them, which keeps the split reproducible from the inventory alone and keeps the
discrepancy countable. This paper reports the discrepancy rather than the larger
file count.

\subsection{The pre-registered subsample}

Labelling was applied to a subsample of \NCohort{} reports from the frozen test
split, which is \CoveragePct{} per cent of it. The subsample is not the first
\NCohort{} identifiers in any ordering that correlates with content: the
selection rule sorts test identifiers by a salted hash and takes the first
\NCohort{}, with the salt fixed and recorded before the labels were read. The
recorded contract states explicitly that the result is not a lexicographic
prefix.

This matters for interpreting one of the results below. A pre-registered
selection rule protects against choosing a sample after seeing its labels. It
does not protect against the sample being unusual, and the two are routinely
conflated.

\subsection{Labels}

Labels are the \NFindings{} observations of the CheXpert schema~\cite{irvin2019chexpert},
assigned to report text by the released CheXbert labeller~\cite{smit2020chexbert},
which is the labeller the standard automatic metric in this area is built on.
The checkpoint used is recorded with its SHA-256 checksum in the recorded artifact, so
the labelling can be tied to a specific set of weights rather than to a name.

The labeller emits four values per observation: positive, negative, uncertain,
and not mentioned. Throughout this paper an observation counts as a gold positive
only when the labeller assigns the positive value. Uncertain and unmentioned are
not counted as positives, and they are not counted as negatives either when a
count of positives is what is being reported.

The unmentioned value dominates and its interpretation is a genuine caveat.
Across the \NFindings{} observations the median count of test reports that do not
mention an observation at all is \MedianBlankPct{} per cent of the split, and for
\MaxBlankFinding{} it is \MaxBlank{} of \NTest{} reports. A report that is silent
about an observation is not a report asserting its absence, and any analysis that
treats the two as equivalent is making a convention visible only in its methods
section. We avoid the convention by reporting positives and misses on positives
rather than accuracy over all four values.

For the four-state composition figure, let \(n_{f,s}\) be the count of state
\(s\in\{+, -, ?, \mathrm{blank}\}\) for finding \(f\) in the frozen test split.
The build checks \(\sum_s n_{f,s}=\NTest{}\) for every finding and plots the
fractions \(n_{f,s}/\NTest{}\). The blank state is therefore retained as a
separate denominator component and is never interpreted as a negative label;
the resulting display describes CheXbert output rather than clinical prevalence.

\subsection{The presence gate and the two power questions}

The analysis applies a presence gate fixed before the runs: an observation is
scored only if the subsample holds at least \MinPresent{} gold positives for it,
and otherwise its cell is reported as underpowered with no rate attached. The
gate value is not asserted here but recovered: the recorded per-observation
status of both control arms is checked against the positive counts, and the build
stops if any status disagrees with a gate of \MinPresent{}.

We keep the eligibility sets explicit. Let \(S\) be the frozen test reports,
\(C\subseteq S\) the \NCohort{}-report draw, and \(P_f=\{i\in S:
y_{if}=+\}\) the reports labelled positive for finding \(f\). The pilot
denominator is
\[
  K_f = |C\cap P_f|,
  \qquad
  \mathrm{status}(f)=
  \begin{cases}
    \mathrm{scored}, & K_f\geq \MinPresent{},\\
    \mathrm{underpowered}, & K_f<\MinPresent{}.
  \end{cases}
\]
At the benchmark level, a cell is eligible only when the permission,
artifact-availability, and no-substitution gates all hold; those gates are
recorded separately from \(K_f\). Thus an inaccessible cell, an observable but
underpowered cell, and a scored cell have different denominators and are never
collapsed into one missing-value code.

Two distinct power questions follow, and separating them is one of this paper's
points.

The first asks what the corpus can support. For an observation present in
\(K\) of the \(N\) test reports, a draw of \(n\) reports contains \(nK/N\)
positives in expectation, so the smallest draw whose expectation reaches the gate
is \(\lceil \text{\MinPresent{}} \cdot N / K \rceil\). When that number exceeds
\(N\) itself, no subsample of the test split can reach the gate and the
limitation is the corpus, not the pilot. This is arithmetic on counts, not a
simulation.

The second asks whether the draw that was actually taken resembles the split it
came from. For each observation we compute the exact two-sided hypergeometric
probability of the observed positive count, summing the probability of every
outcome no more likely than the one observed, under the null that the draw is
uniform over the split. With \NFindings{} observations tested the Bonferroni
threshold is \Bonferroni{}, and we report probabilities against that threshold
rather than against \(0.05\).

\subsection{The two controls}

Two controls are scored in label space on the subsample. The majority control
emits, for every report, the single most frequent findings text in the training
split; it is constant by construction, which the analysis verifies by requiring
each of its per-observation positive counts to be either zero or all \NCohort{}
rows. The retrieval control embeds the indication field with term-frequency--inverse-document-frequency (TF--IDF) weighting, retrieves
the nearest training report by cosine similarity, and returns that report's
findings text; retrieval is restricted to the training side, so a report can
never retrieve itself or a study from its own patient.

For each observation and each control we count the gold positives the control
failed to assert, and report that count as a rate with a one-sided \(90\%\)
Clopper--Pearson upper limit~\cite{clopper1934use}. The upper limit is reported
because the counts are small: at these sample sizes an interval is the result and
a point estimate on its own would misrepresent it.

For a scored finding \(f\), the reported miss rate is
\(\widehat q_f=M_f/K_f\), where \(M_f\) is the number of positive reports the
control failed to assert. The denominator is therefore the observed
gold-positive count in \(C\), not \(|C|\), \(|S|\), or the number of files on
disk. No \(\widehat q_f\) is defined when \(K_f<\MinPresent{}\).

The source design asked for a coverage--risk curve over abstention rates.
The bound control tables record a single coverage of the labelled draw and
no ranking that would order reports for intermediate abstention. Only the
\NCoverageRiskPoints{} cells with \(K_f\geq\MinPresent{}\) can contribute a
point, and each point is \(\widehat q_f\) at coverage
\CoverageRiskCoverage{}. A curve is therefore not drawn. The same tables
record the contact-gated evaluation as \TypeRBound{}; no score is computed
or printed under that name.

The same two controls also exist scored a different way, on surface text over the
whole test split rather than in label space on the subsample, using exact match
and unigram token F1. Those numbers are reported here only as a contrast, and
they are not comparable to the label-space numbers.

\subsection{Quantities excluded by design}

Three quantities were within the project's reach and were kept out of the
analysis by design, because each would present the result as stronger than
the access situation supports.

Scores are reported under the name of the corpus they were computed on. The
locally held corpus is a different corpus from the one the benchmark scores,
and a number computed on it carries that corpus's name rather than the
benchmark's, since the substitution would be a category error rather than an
approximation. The project's own scoring wrapper enforces this by accepting
only samples of the benchmark's split size.

The label counts are reported as counts of positives and misses. They are kept
out of the field's headline automatic report-generation aggregate, because
that conversion would invite comparison against published values computed on a
different corpus with a different sample size.

The one local inference run is reported as what it is. A run of a
chest-radiograph foundation model over \HarnessImages{} images exists and is
recorded as a smoke test that checks the code runs end to end, explicitly
marked as not the official evaluation. It contributes no number to this paper.

\subsection{Evidence discipline}

No number in this manuscript is typed by hand. Every artifact it rests on is
recorded with its SHA-256 checksum in a manifest that accompanies the manuscript
source; the figure and table code verifies each checksum against the bytes on disk
before reading them and aborts on any mismatch, and it emits both the figures and
the numeric macros the text prints. Several consistency checks are enforced at
build time rather than asserted in prose: the two control arms must agree on how
many gold positives the subsample holds, the recomputed list of observations
expected to yield fewer than one positive must match the list the runs recorded,
the presence gate must be recoverable from the recorded statuses, and the set of
access obstacles the first figure describes must equal the set the project
recorded. Any disagreement stops the build instead of producing a manuscript.


\section{The blockage is administrative}

Fig.~\ref{fig:locks} lays out the \NLocks{} recorded obstacles, grouped by
what would remove each one. The grouping is the result, and the three groups
behave differently enough that treating them as one queue would be a mistake.

Of the \NLocks{}, \NPermissionLocks{} are permissions. The benchmark suite is
distributed under a contact gate, so its maintainers grant access on
request~\cite{chexbench}. Human-subjects research training must be completed by
a named person and produces a verification record in that person's name. The
radiographs the benchmark scores are governed by a credentialed data use
agreement of their own~\cite{johnson2019mimiccxr,goldberger2000physionet}. The
project's record is explicit that these three are not interchangeable: holding
any one of them satisfies neither of the others, and none of them is a token
that can be substituted for another. That sounds obvious written down and is
routinely got wrong in planning, because all three are colloquially ``getting
access''.

Fig.~\ref{fig:permissions} draws the relation the record states, with the
permissions on one axis and the things a project needs on the other. Each
approval delivers exactly one column and leaves the other
\NPermissionPairs{} cells untouched, so the grid is a diagonal rather than a
progression. Reading it as a progression is the specific error the shape rules
out: the access request yields the file defining which cases are scored and no
images, the training certificate yields a credential that later agreements
require of a named person and no data, and the data use agreement yields the
radiographs and neither of the other two. A plan that treats the first
acceptance as a third of the way to running the benchmark has mistaken three
independent conditions for one queue, and will report progress it does not
have.

Exactly \NArtifactLocks{} of the \NLocks{} is an artifact rather than a
permission. The file
defining which cases the benchmark scores was checked for at three locations and
found at none of them. It is not a fourth door: it arrives with the first
permission. It is listed separately because holding a permission and holding the
file it unlocks are different states, and a plan that tracks only permissions
will report itself further along than it is.

The record also shows the absence was not for want of looking. Public mirrors
and archive copies were checked and none held the file, and one similarly sized
file circulating under a related name was examined and rejected because it
defines a different task. That refusal is the part worth copying. A file of
about the right size under about the right name is the most likely way an
evaluation quietly becomes a different evaluation, and the only defence is to
check what a candidate file defines before running anything against it.

The remaining \NSubstitutionLocks{} obstacles are refusals the project imposed
on itself, and they are the only ones it controlled. A different public
chest-radiograph corpus was available locally, and scoring it under the
benchmark's name was refused as a category error rather than accepted as an
approximation. The project's scoring wrapper enforces the refusal mechanically
by rejecting any sample whose size is not the benchmark's own split size, which
is why the \HarnessImages{}-image local inference run recorded in the same
project is recorded as a harness that checks the pipeline executes and
contributes no number to this paper.

What makes the grouping worth publishing is a negative property that holds
across all three. None of these obstacles is removed by more computation, by a
larger model, or by labelling more of the corpus already on hand. A project that
reads its own blockage as a resource problem will spend resources and remain
blocked. The record here was written before that spending, and stating that
plainly is the only way the avoided cost shows up anywhere.

The same record sets out what a result would take, and its shape is the
clearest statement of the claim. Table~\ref{tab:unlock} reproduces the route: it
runs to \NUnlockSteps{} steps, of which \NMachineSteps{} are mechanical and
exactly one needs a person. That one is the first, and nothing after it can
start until it finishes. This is why the study reports the blockage rather than
an attempt at the benchmark. A route whose only human step is at the head cannot
be shortened from the other end, and the steps that a rented machine could
perform are precisely the ones that were never reached. Publishing the route
also makes the claim falsifiable in the ordinary way: a reader who holds the
approvals can execute the remaining steps and check whether they behave as the
record says.

\begin{figure}[tbp]
\centering
\includegraphics[width=\linewidth]{fig1_access_locks.pdf}
\caption{\CapLocks}
\label{fig:locks}
\end{figure}

\begin{figure}[tbp]
\centering
\includegraphics[width=\linewidth]{fig2_permissions.pdf}
\caption{\CapPermissions}
\label{fig:permissions}
\end{figure}

\begin{table}[tbp]
\centering
\caption{The route the project recorded from its present state to a result on
the benchmark, written down before any machine was rented. Step is the recorded
order; a step cannot begin until every step above it has finished. Who acts
distinguishes the steps that need a person from those a machine performs
unattended. What it waits on is the condition the record attaches to each step,
restated as the state of the world that must hold rather than as the command
that checks it. The head of the route is the only step a person must take, and
it gates every other.}
\label{tab:unlock}
\small
% Generated by the figure script from hash-bound evidence. Do not hand-edit.
\begingroup
\setlength{\tabcolsep}{5pt}
\begin{tabular}{@{}rp{0.28\linewidth}lp{0.40\linewidth}@{}}
\toprule
Step & Action & Who acts & What it waits on \\
\midrule
0 & Obtain the three approvals & A person & Nothing. No machine is rented and no gated file is fetched before the approvals are held, which is why the cost of being blocked here is zero. \\
1 & Fetch the file defining which cases are scored & Machine & The access request has been granted, and the fetch is made by the rented machine under its own credential. \\
2 & Check the file before spending machine time & Machine & The file parses, matches the task's schema, and holds the published number of cases. A file that merely resembles the right one stops here. \\
3 & Obtain the images the file names & A person, then machine & Each image source's own agreement is held. Which agreements those are is decided by the file, not by the project. \\
4 & Generate a report for every case & Machine & The three preceding steps have passed and the model is already on the machine. \\
5 & Score the reports with the benchmark's own labeller & Machine & The generated reports exist, and the scoring uses the benchmark's labeller rather than an approximation of it. \\
6 & Copy the scores back & Machine & The scores were verified where they were computed. This step moves evidence and computes nothing. \\
7 & Record the result & Machine & The scoring wrapper is given the benchmark's real split size, and its refusal of undersized samples stays active. \\
8 & Close the question & The project & The recorded result is committed and reviewed independently of the person who produced it. \\
\bottomrule
\end{tabular}
\endgroup

\end{table}

\section{The corpus sets a ceiling the pilot cannot raise}

The pilot labelled \NCohort{} of the \NTest{} test reports, \CoveragePct{} per
cent of the split, and \NUnderpowered{} of the \NFindings{} observations fall
below the presence gate. The obvious reading is that the pilot is too small. For
most of the observations that reading is right, and for \NExceedSplit{} of them
it is wrong in a way that matters.

Table~\ref{tab:power} and Fig.~\ref{fig:required} give the arithmetic. For
each observation the table reports how many labelled reports would be needed
before a draw is expected to hold \MinPresent{} positives, given that
observation's prevalence in the split. For \NExceedSplit{} observations that
number exceeds \NTest{}, the size of the whole test split: \RarestFinding{} is
positive in \RarestPresent{} of \NTest{} reports and would need
\RarestRequired{} of them. Labelling every remaining report in the split would
not move those cells above the gate. One observation,
\BoundaryFinding{}, sits exactly at the boundary and would need the split
entire.

The distinction this draws is between a pilot that is underpowered and a corpus
that cannot answer the question. They are different findings with different
consequences. The first is fixed by labelling more of what is already on hand,
which costs annotation time and nothing else. The second is fixed only by a
larger or differently composed corpus, which is a different project. A report
that says only ``underpowered'' points the reader at the cheaper of the two
remedies for cells where the cheaper remedy cannot work.

It is worth being precise about what the required-rows column is not. It is the
size at which the \emph{expected} number of positives reaches the gate, which is
a necessary condition and a weak one: at exactly that size roughly half of all
draws still fall short. Read it as a floor below which the question is
hopeless, not as a sample size recommendation.

\begin{table}[tbp]
\centering
\caption{What the frozen test split of \NTest{} reports contains for each of the
\NFindings{} labelled observations, and what a subsample would need. Rows needed
for the gate is the smallest draw whose expected number of gold positives
reaches \MinPresent{}, computed from that observation's prevalence in the split;
values in bold exceed the size of the split itself, so no subsample of it can
reach the gate. Drawn is the number of gold positives in the \NCohort{}-report
subsample that was actually labelled and Expected is what the split's prevalence
predicts for a draw of that size. Unmentioned counts reports whose text does not
refer to the observation at all, which the labeller distinguishes from
asserting its absence.}
\label{tab:power}
\small
% Generated by the figure script from hash-bound evidence. Do not hand-edit.
\begingroup
\setlength{\tabcolsep}{5pt}
\begin{tabular}{lrrrrrl}
\toprule
Observation & Positive & Unmentioned & Rows needed & Drawn & Expected & Status \\
& in split & in split & for the gate & & & \\
\midrule
Enl. Cardiomediastinum & 5 & 128 &   333 & 0 & 0.5 & underpowered \\
Cardiomegaly & 41 & 65 &    41 & 3 & 3.9 & underpowered \\
Lung Opacity & 44 & 238 &    38 & 12 & 4.2 & scored \\
Lung Lesion & 18 & 299 &    93 & 2 & 1.7 & underpowered \\
Edema & 3 & 307 & \textbf{  555} & 0 & 0.3 & underpowered \\
Consolidation & 1 & 214 & \textbf{1,665} & 0 & 0.1 & underpowered \\
Pneumonia & 4 & 313 & \textbf{  417} & 1 & 0.4 & underpowered \\
Atelectasis & 14 & 312 &   119 & 1 & 1.3 & underpowered \\
Pneumothorax & 1 & 75 & \textbf{1,665} & 0 & 0.1 & underpowered \\
Pleural Effusion & 8 & 74 &   209 & 0 & 0.8 & underpowered \\
Pleural Other & 3 & 330 & \textbf{  555} & 0 & 0.3 & underpowered \\
Fracture & 17 & 298 &    98 & 2 & 1.6 & underpowered \\
Support Devices & 18 & 315 &    93 & 2 & 1.7 & underpowered \\
No Finding & 130 & 203 &    13 & 8 & 12.5 & scored \\
\bottomrule
\end{tabular}
\endgroup

\end{table}

\begin{figure}[tbp]
\centering
\includegraphics[width=\linewidth]{fig3_required_rows.pdf}
\caption{\CapRequired}
\label{fig:required}
\end{figure}

\section{A pre-registered draw is still a draw}

The \NCohort{} reports were selected by sorting test identifiers on a salted
hash and taking the first \NCohort{}, with the salt fixed and recorded before
any label was read. That protocol rules out the failure mode it was designed
for. It does not rule out the draw being unusual, and this one is.

Fig.~\ref{fig:draw} compares, for every observation, the positives the draw
contains against the positives the split's own prevalence predicts for a draw of
that size. Exactly \NUnrepresentative{} of them is out of line: \SkewFinding{} is
positive in \SkewSplitPresent{} of \NTest{} test reports, which predicts
\SkewExpected{} positives in a draw of \NCohort{}, and the draw holds
\SkewObserved{}.
The exact two-sided hypergeometric probability is \SkewP{}, against a Bonferroni
threshold of \Bonferroni{} for \NFindings{} comparisons. Every other observation
is consistent with the split, and \NZeroDrawn{} of them contribute no positives
at all.

The consequence is not cosmetic. Only \NScored{} of \NFindings{} observations
clear the presence gate, and \SkewFinding{} is one of them --- with
\ScoredOnePresent{} positives against \ScoredOneExpected{} expected. Had the
draw been typical, that cell would have fallen below the gate with the rest. At
the expected draw, \NScoredExpected{} of the \NFindings{} observations clears
the gate: \ScoredTwo{}, which is positive in \ScoredTwoSplit{} of \NTest{}
reports and is the only observation common enough to survive a draw this small
reliably. One of the pilot's \NScored{} scorable cells therefore exists because
a pre-registered random draw happened to be lucky, or unlucky, depending on
which way one reads it.

We want to state the implication carefully, because it is easy to overreach.
Nothing here indicates a defect in the selection rule, which did what a salted
hash does. Nothing here indicates the labels are wrong. The observation is
narrower and, we think, more useful: at these sizes the set of cells a pilot can
report on is itself a random variable, and a pilot that reports only its
scorable cells is reporting a selection it did not make and usually does not
mention. The remedy is cheap --- state the expected count next to the observed
one for every cell, scorable or not, which is what the table above does --- and
we have not seen it done.

\begin{figure}[tbp]
\centering
\includegraphics[width=\linewidth]{fig4_draw_vs_split.pdf}
\caption{\CapDraw}
\label{fig:draw}
\end{figure}

\section{Both controls fail, and one of them fails informatively}

\subsection{What the controls assert}

Fig.~\ref{fig:assertions} shows what each control asserted across all
\NFindings{} observations, next to what the draw contains. The retrieval control
behaves as a retrieval control should: it asserts positives for
\RetAssertedFindings{} different observations, in counts that are the right
order of magnitude, because it returns real reports written by radiologists
about other patients.

The majority control does something else. It emits one constant report for every
row, so each of its per-observation counts must be zero or all \NCohort{}, which
the analysis verifies rather than assumes. It reaches \NCohort{} for exactly one
observation, \MajAsserted{}, and zero for the other thirteen. The text it emits
is the most frequent findings text in the training split:

\begin{quote}
\small\itshape \MajTemplate{}
\end{quote}

That reads as an unremarkable study. The labeller marks it positive for
\MajAsserted{} on every row and never marks it as \ScoredTwo{}, the observation
that denotes a study with nothing to report. So the majority control asserts
\MajAsserted{} on all \NCohort{} reports, where the truth is
\MajAssertedDrawn{}, and asserts nothing at all on the two cells that can be
scored.

This is worth separating from the pilot it was found in. A constant majority
baseline is standard, and its purpose is to establish the floor that a system
predicting nothing informative would achieve. In label space that purpose fails
here, and not because the corpus is small: it fails because the floor is defined
by whatever the labeller assigns to one particular string, and that assignment
is not the null vector anyone assumed. Any comparison in which the majority
baseline's label-space score is treated as the trivial floor inherits this,
whatever its sample size.

\subsection{The two cells that can be scored}

Table~\ref{tab:scored} and Fig.~\ref{fig:scored} give the two cells that clear
the gate. On \ScoredOne{}, with \ScoredOnePresent{} gold positives, the majority
control misses all of them and the retrieval control misses \RetMissedOne{}, a
share of \RetMissOne{} with a one-sided \(90\%\) upper limit of \RetUpperOne{}.
On \ScoredTwo{}, with \ScoredTwoPresent{} gold positives, the majority control
again misses all of them and the retrieval control misses \RetMissedTwo{}, a
share of \RetMissTwo{} with an upper limit of \RetUpperTwo{}.

The retrieval control is better than the majority control in both cells, and the
correct summary of that is not that retrieval works. Both upper limits sit close
enough to one that these data cannot exclude a control that misses essentially
every positive. Two cells with \ScoredOnePresent{} and \ScoredTwoPresent{}
positives do not distinguish a control that misses four fifths of positives from
one that misses all of them, and reporting the point estimates without the
limits would suggest otherwise.

\subsection{The same controls, scored differently}

The two controls also exist scored on surface text over the whole \NTest{}-report
split, by exact match and unigram token F1. There, retrieval reaches
\SurfaceRetrievalF{} token F1 on the findings section against \SurfaceMajorityF{}
for the majority control, a ratio of about \SurfaceRatio{} to one, and exact
match is \SurfaceRetrievalExact{} per cent against \SurfaceMajorityExact{} per
cent.

Those numbers are calm and orderly and they describe the same two systems that,
in label space, respectively fail to exclude near-total failure and assert one
wrong observation on every report. Surface overlap rewards a control for
producing text of the right shape and vocabulary, which both controls do by
construction, since both emit real report text. Neither number is wrong. They
are measuring the property that survives when the clinical content is removed,
and the gap between the two views of the same pair of systems is the practical
argument for label-space evaluation being worth its cost.

\begin{figure}[tbp]
\centering
\includegraphics[width=\linewidth]{fig5_control_assertions.pdf}
\caption{\CapAssertions}
\label{fig:assertions}
\end{figure}

\begin{table}[tbp]
\centering
\caption{The \NScored{} observations that clear the presence gate of
\MinPresent{} gold positives in the \NCohort{}-report subsample, and how each
control performed on them. Missed counts the gold positives the control did not
assert, miss rate is that count over the gold positives, and the final column is
the one-sided \(90\%\) Clopper--Pearson upper limit on the miss rate, which is
wide because the counts are small. The majority control emits one constant
report for every case; the retrieval control returns the training report whose
indication text is nearest by TF--IDF cosine similarity.}
\label{tab:scored}
\small
% Generated by the figure script from hash-bound evidence. Do not hand-edit.
\begin{tabular}{llrrrr}
\toprule
Observation & Control & Gold positives & Missed & Miss rate & 90\% upper \\
\midrule
Lung Opacity & Majority template & 12 & 12 & 1.00 & 1.00 \\
 & Indication retrieval & 12 & 10 & 0.83 & 0.95 \\
No Finding & Majority template & 8 & 8 & 1.00 & 1.00 \\
 & Indication retrieval & 8 & 7 & 0.88 & 0.99 \\
\bottomrule
\end{tabular}

\end{table}

\begin{figure}[tbp]
\centering
\includegraphics[width=0.86\linewidth]{fig6_scored_cells.pdf}
\caption{\CapScored}
\label{fig:scored}
\end{figure}

\section{Discussion}

The obstacles in this paper are separable, and each is owned by someone
different. The administrative one is owned by institutions and by the person
willing to complete the paperwork, and no amount of engineering substitutes. The
corpus ceiling is owned by whoever assembles corpora, and it is invisible to a
project that only reports its own sample size. The composition of the draw is
owned by nobody, which is why it needs to be printed rather than assigned. The
instrument behaviour is owned by the evaluation community, and it is the only
one a reader of a results table might otherwise mistake for a property of the
system being evaluated.

Reporting them separately costs almost nothing. Everything in this paper is
arithmetic over label counts that already existed, plus one exact test, plus the
decision to publish a list of access obstacles instead of a sentence saying the
benchmark was unavailable. What it buys is that each part can be refuted
independently. Someone who obtains the access grant refutes the first part.
Someone who assembles a corpus with more positives for the rare observations
refutes the corpus ceiling. Someone who shows that the labeller reads the
constant template differently under a different checkpoint refutes the last, and
we have recorded the checkpoint digest precisely so that can be attempted.

Two points generalise past this corpus and this benchmark.

The first is that a pre-registered draw is still a draw. Pre-registration is now
routine and its guarantee is often stated too broadly. Fixing the selection rule
in advance removes selection bias from the choice of sample; it leaves sampling
variation entirely intact, and at pilot sizes sampling variation determines
which cells are reportable at all. The cheap defence is to print the expected
count beside the observed count for every cell, so a reader can see that the
reportable set was drawn rather than designed.

The second is that a trivial baseline is only trivial in the space it is scored
in. The constant majority report used here is the standard construction, and in
surface-overlap space it does what it is supposed to. Passed through an
automatic labeller, it becomes a system that confidently asserts one wrong
observation on every case, which is not a floor in any useful sense. Anyone
reporting a label-space metric with a majority baseline should check what the
labeller assigns to their constant string, and report that vector, because it
takes one forward pass and it determines what the floor means.

\section{What this study does not claim}

No result is reported on the blocked benchmark, and no number appearing here is
presented as an approximation of one. The locally held corpus is a different
corpus and is treated as such throughout.

No standard automatic report-generation aggregate is computed. The counts here
support statements about positives and misses on a small labelled subsample;
they are not converted into a headline metric, and no comparison against
published values is made or implied.

No deficiency is alleged in the benchmark, in its maintainers, or in the
credentialing bodies whose approvals are enumerated here. Gating clinical data
behind institutional agreements is a legitimate and, for these data, an expected
choice. What is reported is the consequence for a project that has not completed
those steps, which is a fact about the project.

No claim is made about the labeller's correctness. That the labeller assigns
\MajAsserted{} to one constant string is an observation about that string under
that checkpoint. It is not evidence that the labeller is inaccurate in general,
and this study did not evaluate its accuracy.

No conclusion is drawn about report-generation systems. Both arms examined here
are controls. Nothing in this paper compares a trained generator against
anything, and the \HarnessImages{}-image local inference run that exists in the
same project is recorded as a pipeline check and contributes no number.

No cell below the presence gate is given a rate, an interval, or a direction.
Of the \NFindings{} observations, \NUnderpowered{} are reported as underpowered
and are left there.

\section{Conclusion}

One chest-radiograph reporting project that did not reach its intended
evaluation was decomposed into the reasons it did not. A set of \NLocks{} access
obstacles was enumerated and grouped by what removes them, and none of the
three groups is removed by computation; the enumeration was recorded before
hardware was rented rather than after. Of \NFindings{} labelled observations,
\NExceedSplit{} require more reports than the entire \NTest{}-report test split
before a draw is expected to hold \MinPresent{} positives, so for those the
limitation is the corpus and not the pilot. The pre-registered \NCohort{}-report
subsample is not representative of that split, over-representing \SkewFinding{}
at \SkewObserved{} positives against \SkewExpected{} expected with an exact
two-sided hypergeometric probability of \SkewP{}, and one of the \NScored{}
cells that clear the presence gate exists only because of that departure.

On both scorable cells the two controls fail, and the one that fails completely
does so for a reason that travels: the constant report a majority baseline emits
is read by the automatic labeller as positive for \MajAsserted{} on every case
and as never denoting an unremarkable study, so the floor that baseline is meant
to establish is not the floor anyone assumed. Scored on surface text the same
pair looks unremarkable at \SurfaceRetrievalF{} and \SurfaceMajorityF{} token
F1. Across every one of these results the pattern is the same: the informative
quantity is the specific reason a claim cannot be made, and every aggregate that
a project like this would normally report --- feasible or not, powered or not,
better or worse --- hides exactly that.

\appendix

\section{The bound evidence}
\label{app:evidence}

The methods section asserts that no number in this manuscript is typed by hand
and that every artifact behind one is pinned by digest. Table~\ref{tab:evidence}
is that assertion in a form a reader can check. It lists all \NEvidence{} bound
artifacts, totalling \EvidenceBytes{} bytes, in the order the analysis consumes
them.

Two things about this table are worth stating, because a manifest printed
without them invites more confidence than it earns. First, a digest binds bytes
to a claim; it does not make the bytes correct. What Table~\ref{tab:evidence}
rules out is drift --- an artifact quietly changing after a number was computed
from it, or a figure being redrawn from a file that is no longer the one the
prose describes. It does not rule out an artifact having been wrong when it was
written, and nothing in this appendix should be read as claiming otherwise.
Second, the internal identifier and the path of each artifact are omitted on
purpose. They describe a private tree, they would mean nothing to a reader
outside it, and keeping them out of the PDF is a property the build checks
rather than a courtesy. What remains is what a reader can act on.

The digests are also what makes the negative results in this paper durable. A
study whose main finding is that something could not be measured is unusually
easy to revise after the fact, because there is no positive number anchoring it.
Recording the obstacle enumeration and the blocked tier under digest, before any
attempt was made, is what stops the account of why this pilot could not answer
its question from being rewritten once the answer is known.

\begin{table}[htbp]
\centering
\caption{Every artifact this manuscript rests on, with its size and the first 16
hexadecimal characters of its SHA-256 digest. Full digests, and the paths the
digests apply to, are in the manifest that accompanies the manuscript source;
the prefix is shown here only because a 64-character digest does not set in a
table column. The figure and table code recomputes each digest from the bytes on
disk before reading them and stops the build on any mismatch, so a stale
artifact cannot reach the page.}
\label{tab:evidence}
\footnotesize
% Generated by the figure script from hash-bound evidence. Do not hand-edit.
\begingroup
\setlength{\tabcolsep}{5pt}
\begin{tabular}{@{}p{0.56\linewidth}rl@{}}
\toprule
What it records & Bytes & SHA-256 (first 16) \\
\midrule
the four access locks on the contact-gated benchmark, each with what it requires and what was checked &   3,751 & \texttt{4b1a1e3e4f777383} \\
the recorded blocker: contact gate plus an absent official split file &     392 & \texttt{9e7ff622a2f6ba52} \\
the full-replication tier, recorded as blocked rather than attempted &     546 & \texttt{1dd3365c638955fc} \\
the recorded finding that the three access requirements are not interchangeable &  11,023 & \texttt{4a8dd1e543aab01d} \\
the frozen patient-disjoint split: hash rule, patient counts, and the images left outside the inventory & 105,795 & \texttt{8550a81f9285eb3b} \\
the pre-registered salted-hash rule that selected the 32 held-out reports &   1,340 & \texttt{4b67f58384bc18ed} \\
per-finding gold label counts over the full frozen test split &   3,269 & \texttt{b3128f550d0be2e4} \\
per-finding present, predicted and missed counts for both control arms on the held-out subsample &   6,236 & \texttt{1d64a6eeb89722d0} \\
the retrieval control scored on surface text over the full test split &     816 & \texttt{e7426042fcb4d85d} \\
the majority control scored on surface text over the full test split &     896 & \texttt{21c4b64cda7cfb90} \\
the eight-image inference run, recorded as a harness rather than a benchmark result &     468 & \texttt{b8dd79a6dc313cc2} \\
the same run entered in the comparison ladder, with its claim kind recorded &     548 & \texttt{85c7bbd6cb5bd9a6} \\
a provenance record of what is on disk, what its licence permits and what was never run &   5,454 & \texttt{279277922081b4b2} \\
labeller battery on the majority template and its first-sentence and placeholder variants &  96,357 & \texttt{18f5f05bf3908de4} \\
the same two control arms scored on the full frozen test split of 333 reports & 232,370 & \texttt{644872c3a1d33f3b} \\
gold labels recomputed after stripping the de-identification placeholder from each test report & 167,285 & \texttt{530c949d97f8436e} \\
execution receipt for the 2026-09-07 labeller battery, 333-row controls, and gold-placeholder ablation &   5,280 & \texttt{40664079cc4d4a26} \\
\bottomrule
\end{tabular}
\endgroup

\end{table}

\section*{Data and code availability}

The derived label tables, the recorded access-obstacle enumeration, the split
and selection protocols, and the figure code that produces every number in this
manuscript accompany the manuscript source. Each artifact is recorded with its
SHA-256 digest in a manifest, and the figure code verifies those digests against
the bytes on disk before reading them, so the build fails rather than silently
reporting stale numbers. The radiograph pixels and the report texts are not
redistributed here; the corpus is publicly available from its original
distributor under its own licence, and the labeller checkpoint is a third-party
release identified in the manifest by digest rather than copied.
%% Generated by scripts/release_targets.py from the release ledger.
%% Identifiers below are ledger-checked. The sentence is a Task-4 honesty
%% pass: a deposit row is not a finished public scholarly release.
%% Ledger state: github=CREATED, zenodo=DEPOSITED
\ifdefined\ANONYMOUS That material is intended for an archive under a persistent identifier; both addresses are withheld from this journal review copy and appear in the identified version. This is not yet a public scholarly release.\else That material is intended for archive under the persistent identifier \href{https://doi.org/10.5281/zenodo.22647018}{10.5281/zenodo.22647018}; a development repository is listed at \url{https://github.com/PeterPonyu/blocked-pilot-decomposition}. This is not yet a public scholarly release.\fi


\section*{Ethics statement}

This study is a secondary analysis of a de-identified, publicly distributed
collection of radiology reports and of derived label counts computed from them.
No new data were collected, no images or report texts are reproduced beyond a
single template sentence that is the modal training-split text and contains no
patient-specific content, and no attempt was made to re-identify any individual.
The credentialed data governed by an agreement that is discussed in this paper
was never obtained and never used.

\section*{Competing interests}

Not yet declared. This is a working draft and the declaration is outstanding.

\section*{Funding and author contributions}

Not yet declared. This is a working draft; the author list is not fixed and
contributions are therefore not yet attributable.

\section*{Declaration of generative AI use}

The analysis code, the figure code and the drafting of this manuscript were
produced with substantial assistance from large language model agents working
under human direction. One safeguard bears directly on what is reported: every
number printed above is emitted programmatically from artifacts bound by
cryptographic digest rather than transcribed by a model into prose, and several
of the paper's structural claims --- that the majority control is constant, that
the presence gate is the value stated, that the recorded and recomputed lists of
thin observations agree --- are enforced as build-time checks that stop the
build rather than as sentences. Responsibility for the content rests with the
authors.
